\documentclass[11pt,a4paper]{article}

\usepackage[T1]{fontenc}
\usepackage[utf8]{inputenc}
\usepackage[french]{babel}
\usepackage{geometry}
\usepackage{amsmath,amssymb}
\usepackage{array}
\usepackage{booktabs}
\usepackage{longtable}
\usepackage{xcolor}
\usepackage{enumitem}
\usepackage{tcolorbox}
\usepackage{parskip}
\usepackage{hyperref}

\geometry{top=2.5cm, bottom=2.5cm, left=3cm, right=2.5cm}

\definecolor{bleuTitre}{RGB}{31,56,100}
\definecolor{bleuSection}{RGB}{31,92,153}
\definecolor{bleuClair}{RGB}{220,235,250}
\definecolor{vertOk}{RGB}{30,120,50}
\definecolor{rougeAlert}{RGB}{180,40,20}
\definecolor{jauneAttention}{RGB}{180,140,30}
\definecolor{grisNote}{RGB}{245,244,240}

\hypersetup{
  colorlinks=true,
  linkcolor=bleuSection,
  urlcolor=bleuSection
}

\newtcolorbox{encadre}[1]{
  colback=bleuClair, colframe=bleuSection,
  title=#1, fonttitle=\bfseries
}

\newtcolorbox{hypothese}{
  colback=grisNote, colframe=jauneAttention,
  title={\textbf{Hypothèse}}, fonttitle=\bfseries
}

\newtcolorbox{limite}{
  colback=grisNote, colframe=rougeAlert,
  title={\textbf{Limite / critique}}, fonttitle=\bfseries
}

\title{\color{bleuTitre}\textbf{Analyse des résultats multi-runs\\[0.2cm]
Pipeline CSP $\to$ xTB sur benzénoïdes h3, h4, h5}}
\author{Stage M2 IMD --- Aix-Marseille Université}
\date{\today}

\begin{document}

\maketitle

\tableofcontents
\bigskip

% ======================================================
\section{Contexte et objectif de l'analyse}
% ======================================================

Le pipeline complet est~:
\[
\text{.graph (BenzAI)} \;\to\; \text{Solveur CSP (ACE)} \;\to\; \text{Reconstruction 3D} \;\to\; \text{xTB GFN2 \texttt{--opt tight}} \;\to\; \text{Test ACP}
\]

\paragraph{Pourquoi multi-runs ?}
xTB est un optimiseur \emph{local}~: il descend dans le puits d'énergie le plus proche de la géométrie initiale. Si la molécule possède plusieurs minima d'énergie (par exemple un conformère plat et un conformère plié), un seul run xTB peut tomber dans l'un ou dans l'autre selon la perturbation initiale ($\pm 0{,}1$~Å en $z$). On lance donc $N=10$ runs par solution avec des perturbations aléatoires différentes pour caractériser la \emph{stabilité} de la classification.

\paragraph{Données analysées.}
24 fichiers \texttt{data.json} (3 valeurs de $h$ $\times$ 8 configurations CSP), 692 solutions au total, 10~runs chacune --- soit \textbf{environ 6920 optimisations xTB exploitées}. Tous les originaux tout-hexagonal sont multi-runs aussi (200 cas, test de cohérence du pipeline).

\paragraph{Classifications appliquées (seuils fixes).}
\begin{itemize}[leftmargin=2em,itemsep=0.2em]
  \item \texttt{always\_planar}~: 100\% des runs plans \emph{et} $\mu_{\text{angle}} < 5\degree$.
  \item \texttt{mostly\_planar}~: $\geq 70\%$ plans \emph{et} $\sigma < 3\degree$.
  \item \texttt{unstable}~: 30--70\% plans \emph{ou} $\sigma > 5\degree$.
  \item \texttt{mostly\_non\_planar}~: $< 30\%$ plans (mais $> 0\%$).
  \item \texttt{always\_non\_planar}~: 0\% plans.
  \item \texttt{ambiguous}~: autres cas.
\end{itemize}

% ======================================================
\section{Tests de cohérence du pipeline}
% ======================================================

\begin{encadre}{Résultat global}
\textbf{Le pipeline est strictement cohérent. Aucune incohérence détectée.}
\begin{itemize}[leftmargin=2em,itemsep=0.2em]
  \item \textbf{200 originaux tout-hexagonal} (= benzénoïdes BenzAI sans substitution)~: 100\% planar, angle maximum 0,87\degree.
  \item \textbf{200 solutions CSP tout-6} (= le solveur reproduit la signature originale)~: 100\% \texttt{always\_planar}.
\end{itemize}
\end{encadre}

\paragraph{Pourquoi c'est important.}
Si une solution tout-6 produite par CSP n'était pas \texttt{always\_planar}, cela signifierait soit que la reconstruction 3D introduit des distorsions artificielles, soit que xTB ne reconnaît pas un benzénoïde plat. Le fait que tous les 200 cas soient parfaits valide les deux étapes en aval du solveur (reconstruction + xTB) avant qu'on commence à interpréter les substitutions 5/7. Toute non-planarité observée par la suite est donc \emph{réellement attribuable} aux pentagones et heptagones.

% ======================================================
\section{Découverte centrale : la bimodalité des unstable}
% ======================================================

\subsection{Ce qu'on observe}

Sur les 91 solutions classées \texttt{unstable} (h4 + h5), \textbf{toutes les 91 ($100\%$)} présentent une distribution des 10~angles ACP en \emph{deux groupes nettement séparés}, avec un écart vide entre les deux d'au moins 5\degree. Aucune n'est uniformément distribuée autour du seuil de 10\degree.

\textbf{Exemples concrets} (les 10 angles d'une solution unstable)~:
\begin{itemize}[leftmargin=2em,itemsep=0.2em]
  \item \texttt{h5 / 2-5-6-7-13 [v0=6 v1=7 v2=7 v3=5 v4=5]}~:
    angles plans = $\{0{,}1\degree;\ 0{,}3\degree\}$,
    angles courbés = $\{30{,}8\degree;\ 30{,}8\degree;\ 30{,}8\degree;\ 30{,}9\degree;\ 30{,}9\degree;\ 30{,}9\degree;\ 31{,}2\degree;\ 31{,}5\degree\}$.
    Écart vide = 30,5\degree.
  \item \texttt{h5 / 2-5-6-7-8 [v0=7 v1=7 v2=6 v3=5 v4=5]}~:
    8 runs aux alentours de 0,3\degree, 2 runs aux alentours de 29\degree.
    Écart vide = 28\degree.
\end{itemize}

\subsection{Comment l'interpréter (ce qui change)}

\paragraph{Ce qu'on aurait pu croire.}
Le mot \og unstable\fg{} évoque naturellement un cas \og incertain numériquement \fg{}~: on imaginerait une molécule dont l'angle \emph{vrai} est aux alentours de 10\degree, et le bruit numérique fait qu'un run la classe à 9\degree (plan) et un autre à 11\degree (non-plan). Augmenter $N$ corrigerait alors l'imprécision.

\paragraph{Ce qui se passe en réalité.}
La bimodalité massive révèle un phénomène \emph{physique}, pas numérique~: la molécule possède \textbf{deux conformères stables distincts}~:
\begin{itemize}[leftmargin=2em,itemsep=0.2em]
  \item un conformère \emph{plat} (angle ACP $\approx 0\degree$), à plus basse énergie ou comparable, qu'on attend pour les molécules aromatiques planes,
  \item un conformère \emph{plié / courbé} (angle ACP $\approx 25\text{--}35\degree$), où la molécule s'organise en une surface non plane (en forme de bol ou de selle), correspondant à un autre minimum local.
\end{itemize}
Selon la perturbation $z$ initiale, l'optimiseur xTB descend dans l'un ou l'autre puits.

\begin{encadre}{Vocabulaire \og minimum local \fg{}}
xTB minimise l'énergie de la molécule. La fonction d'énergie a en général plusieurs \og puits \fg{} (= minima locaux). xTB n'est pas garanti de trouver le \emph{global}~: il s'arrête au premier puits qu'il rencontre depuis le point de départ. Si on part d'une géométrie légèrement différente (perturbation $\pm 0{,}1$~Å en $z$), on peut atterrir dans un puits différent.
\end{encadre}

\subsection{Conséquences pratiques}

\begin{enumerate}[leftmargin=2em,itemsep=0.3em]
  \item L'\texttt{angle\_mean} ($\approx 10\text{--}15\degree$) d'une solution \texttt{unstable} \textbf{n'a aucun sens physique}. Ce n'est ni l'angle d'un conformère, ni une moyenne sur un mouvement~: c'est la moyenne arithmétique de deux populations qui vivent dans des bassins différents. Ne pas l'interpréter comme \og la molécule est à 12\degree en moyenne \fg{}.
  \item Augmenter $N$ de 10 à 30 \textbf{ne stabilisera pas la classification}. Plus de runs donnera juste une meilleure estimation de \emph{la probabilité d'atterrir dans le bassin plan} (par exemple, on saurait que c'est 20\% plutôt que 30\%), mais le caractère bimodal lui-même ne disparaîtra pas. La classification \texttt{unstable} restera \texttt{unstable}.
  \item La classe \texttt{unstable} mériterait d'être renommée \textbf{\og bistable \fg{}} et reportée comme deux nombres distincts~: l'angle moyen du bassin plan, et l'angle moyen du bassin courbé. Ces deux conformères sont des objets chimiques réels et distincts qui devraient être étudiés séparément.
\end{enumerate}

% ======================================================
\section{Pattern \og palindrome de signature \fg{}}
% ======================================================

\subsection{Définition}

Une signature \texttt{sizes} comme $(v_0, v_1, v_2, v_3, v_4) = (5, 7, 7, 5, 6)$ est dite \emph{palindrome} si elle se lit identiquement de gauche à droite et de droite à gauche. Ainsi~:
\begin{itemize}[leftmargin=2em,itemsep=0.2em]
  \item Palindromes~: $(5,7,7,5)$, $(6,5,7,5,6)$, $(5,6,6,6,5)$.
  \item Non-palindromes~: $(5,6,6,7)$, $(5,7,5,7,6)$, $(7,7,6,5,5)$.
\end{itemize}

\subsection{Résultat sur h5 (527 solutions)}

\begin{center}
\begin{tabular}{@{}lrrrr@{}}
\toprule
 & \textbf{$n$} & \texttt{always\_planar} & \texttt{unstable} & \texttt{non\_planar*} \\
\midrule
Palindromes      & 153 & \textbf{94\%} & 4\%  & 2\%  \\
Non-palindromes  & 374 & 58\%          & 20\% & 22\% \\
\bottomrule
\end{tabular}
\\[0.2cm]
{\small *agrège \texttt{mostly\_non\_planar} et \texttt{always\_non\_planar}.}
\end{center}

\paragraph{Lecture.} Une signature palindrome a $94\%$ de chances d'être \texttt{always\_planar}, contre $58\%$ pour une signature non-palindrome. C'est un \textbf{facteur 1,6 sur la probabilité d'être plat}, et un facteur 11 sur la probabilité d'être non-plan ($2\%$ vs $22\%$).

\begin{hypothese}
La symétrie miroir de la séquence des tailles de cycles correspond (au moins partiellement) à une symétrie géométrique réelle de la molécule. Quand le pentagone à la position $v_0$ a un \og jumeau \fg{} (autre pentagone à la position $v_n$ symétrique), les défauts angulaires qu'ils introduisent se compensent par symétrie axiale. La molécule reste plane parce que chaque déformation à un endroit est équilibrée par sa déformation miroir.
\end{hypothese}

\begin{limite}
Le mot \og palindrome \fg{} ici porte sur la \emph{séquence} $v_0, v_1, \ldots, v_n$, pas sur la molécule elle-même. L'ordre des $v_i$ dépend de l'ordre dans lequel les hexagones apparaissent dans le fichier \texttt{.graph}, qui n'est \emph{pas canonisé} pour la symétrie de la molécule.

Donc~:
\begin{itemize}[leftmargin=2em,itemsep=0.2em]
  \item Une signature palindrome correspond \emph{généralement} à une symétrie géométrique, mais pas toujours.
  \item Une signature non-palindrome peut être en réalité géométriquement symétrique si la numérotation des hexagones ne respecte pas l'axe de symétrie de la molécule.
\end{itemize}
Cela veut dire que le $94\%$ est un \textbf{plancher}~: la vraie corrélation symétrie/planarité est probablement encore plus forte. Pour faire propre, il faudrait recanoniser les indices des hexagones par les automorphismes du graphe avant de tester le palindrome.
\end{limite}

% ======================================================
\section{Pattern \og adjacence 5-7 \fg{} (compensation locale)}
% ======================================================

\subsection{Idée géométrique}

Les angles intérieurs des polygones réguliers sont~:
\[
\alpha(n) = \frac{(n-2) \cdot 180\degree}{n}
\quad \Rightarrow \quad
\alpha(5) = 108\degree, \quad \alpha(6) = 120\degree, \quad \alpha(7) \approx 128{,}57\degree.
\]
Par rapport à un hexagone régulier (référence d'un pavage plan parfait), un pentagone introduit un \emph{défaut angulaire} de $-12\degree$ par sommet (la chaîne se concave, comme dans un fullerène), et un heptagone introduit un défaut de $+8{,}57\degree$ par sommet (la chaîne s'évase, comme dans un graphène hyperbolique).

Quand un pentagone et un heptagone sont \textbf{adjacents} (ils partagent une arête), leurs défauts agissent localement et de signe opposé~: la déformation de l'un se compense partiellement par celle de l'autre sur l'arête commune. Le motif chimique de référence est l'\textbf{azulène} (pentagone fusionné à un heptagone)~: 10 atomes de carbone, structure plate stable, parfaitement aromatique.

\subsection{Mesure quantitative}

J'ai reconstruit le \emph{graphe dual} hex--hex pour chaque molécule (un sommet par hexagone, une arête entre deux hexagones partageant une arête de la molécule), puis pour chaque solution j'ai compté les paires d'hexagones adjacents par taille.

\begin{encadre}{Vocabulaire \og graphe dual \fg{}}
La molécule est un graphe d'atomes (sommets = carbones, arêtes = liaisons C--C). Le \emph{graphe dual} a un sommet par cycle (hexagone, pentagone, heptagone) et une arête entre deux cycles qui partagent une liaison C--C. C'est le bon objet pour parler de \og hexagones voisins \fg{} sans ambiguïté.
\end{encadre}

\textbf{Résultats sur h5} (399 solutions avec au moins une substitution 5 ou 7)~:

\begin{center}
\begin{tabular}{@{}lrrr@{}}
\toprule
\textbf{Sous-groupe} & \textbf{$n$} & \textbf{\% AP} & \textbf{\% non-plan} \\
\midrule
Tous les 5 et 7 ont un partenaire opposé voisin    & 222 & \textbf{73\%} & 14\% \\
Au moins un 5 ou 7 est isolé (pas de partenaire opposé voisin) & 177 & 41\%          & 24\% \\
Que des 5--5 ou 7--7 adjacents (sans 5--7)         &   7 & \textbf{0\%}  & 43\% (+57\% bistable) \\
\bottomrule
\end{tabular}
\end{center}

\textbf{Mêmes résultats sur h4} (79 solutions non-pures)~:
\begin{itemize}[leftmargin=2em,itemsep=0.2em]
  \item Tous appariés~: \textbf{84\% AP}.
  \item Au moins un isolé~: \textbf{48\% AP}.
\end{itemize}

\paragraph{Lecture.}
Passer d'une configuration où chaque pentagone a un heptagone voisin (et vice-versa) à une configuration où au moins un défaut est isolé fait chuter la planarité de \textbf{73\% à 41\%} en h5, et de \textbf{84\% à 48\%} en h4. C'est un effet \emph{énorme} (facteur ${\sim}1{,}8$), reproduit sur deux tailles indépendamment.

\begin{hypothese}
La planarité d'une molécule polycyclique substituée se construit \emph{localement}~: chaque paire 5--7 adjacente forme une sous-unité azulène-like dont les défauts angulaires se compensent à l'échelle d'une arête. Quand un pentagone est isolé (pas de 7 voisin pour le compenser), son défaut $-12\degree$ par sommet \emph{diffuse} sur les hexagones voisins et finit par plier la molécule hors-plan.
\end{hypothese}

\paragraph{Connexion avec le solveur CSP.}
La contrainte \texttt{--adj-57} force chaque pentagone à avoir au moins un heptagone voisin (et vice-versa). \textbf{C'est exactement la condition d'appariement local} de notre hypothèse. C'est pourquoi cette config seule produit \textbf{100\% \texttt{always\_planar} sur h3, h4 et h5} (mais avec une couverture très réduite, voir~\S\ref{sec:configs}). \texttt{adj-57} est essentiellement une formalisation du \og prior \fg{} géométrique \og 5 et 7 doivent être voisins \fg{}.

% ======================================================
\section{Pattern \og topologie de la molécule parente \fg{}}\label{sec:topo}
% ======================================================

\subsection{Définition}

J'appelle \emph{topologie parente} la forme du graphe dual hex--hex de la molécule \emph{avant} substitution. Trois grandes classes apparaissent~:
\begin{itemize}[leftmargin=2em,itemsep=0.2em]
  \item \textbf{Compacte}~: au moins un hexagone a 3 voisins ou plus dans le graphe dual (motif pyrène, triphénylène, coronène).
  \item \textbf{Linéaire}~: tous les hexagones sont alignés sur un même axe (motif anthracène, tétracène, pentacène).
  \item \textbf{Zigzag}~: chaîne d'hexagones sans branchement, mais qui change de direction (motif phénanthrène, chrysène).
\end{itemize}

\subsection{Résultat par molécule h4}

Toutes les solutions substituées (= avec au moins un 5 ou un 7), toutes configs CSP confondues~:

\begin{center}
\begin{tabular}{@{}lcrrr@{}}
\toprule
\textbf{Molécule} & \textbf{Topologie} & \textbf{Solutions substituées} & \textbf{\% AP} \\
\midrule
\texttt{0-5-6-11}  & compacte (pyrène-like)   & 34/34 & \textbf{100\%} \\
\texttt{1-5-6-11}  & compacte                 & 14/15 & \textbf{93\%}  \\
\texttt{0-6-12-18} & chaîne linéaire (tétracène) & 8/8   & \textbf{100\%} \\
\texttt{0-6-11-17} & chaîne zigzag           &  1/8  & \textbf{12\%}  \\
\texttt{0-6-12-17} & chaîne zigzag           &  1/8  & \textbf{12\%}  \\
\texttt{1-5-6-12}  & branchée \og L \fg{}     &  0/6  & \textbf{0\%}   \\
\bottomrule
\end{tabular}
\end{center}

\subsection{Trois régimes nets}

\begin{enumerate}[leftmargin=2em,itemsep=0.3em]
  \item \textbf{Compactes ou linéaires} ($\geq 90\%$ AP)~: la topologie tolère les substitutions 5/7. Le pyrène distribue les défauts sur les 4 directions autour de son sommet central. Le tétracène a un axe long unique sur lequel un pentagone à une extrémité et un heptagone à l'autre s'annulent.
  \item \textbf{Zigzags} (${\sim}12\%$ AP)~: la chaîne change de direction à chaque coude. Les défauts angulaires d'un 5 et d'un 7 ne sont plus alignés sur un même axe et ne se compensent pas géométriquement. Chaque coude introduit une contrainte qui se cumule.
  \item \textbf{Branchée asymétrique} (\texttt{1-5-6-12}, $0\%$ AP)~: structure en \og L \fg{}. Aucune symétrie disponible pour compenser. Pire cas observé.
\end{enumerate}

\begin{hypothese}
Une molécule polycyclique tolère les substitutions 5/7 quand sa topologie offre un \textbf{axe de symétrie} (linéaire) ou un \textbf{centre de symétrie 2D} (compacte type pyrène). Les zigzags et les structures branchées asymétriques ne disposent pas de tels éléments de symétrie naturels~: les défauts angulaires des 5/7 ne se compensent pas et la molécule plie.
\end{hypothese}

\begin{limite}
6 molécules h4, c'est très peu pour parler d'un pattern statistiquement robuste. En particulier le \og 0\% \fg{} sur \texttt{1-5-6-12} repose sur 6 solutions seulement. Il faut tester d'autres molécules de chaque type avant de conclure. C'est précisément le pattern à valider en priorité sur h6 (où 68 molécules sont déjà disponibles).

Note aussi~: il y a probablement un \textbf{couplage} entre l'effet \og adjacence 5-7 \fg{} et l'effet \og topologie compacte \fg{}~: les molécules compactes ont plus d'arêtes hex--hex, donc plus de chances que CSP place les 5 et 7 en position adjacente. Une régression jointe serait nécessaire pour démêler les deux effets.
\end{limite}

% ======================================================
\section{Single-run vs multi-run sur h5}
% ======================================================

\subsection{Comparaison}

Tous les \texttt{data.json} h5 sont déjà en multi-run. Le champ \texttt{planar} et \texttt{angle\_deg} en haut d'une solution correspond simplement au \emph{dernier} des 10 runs (rétrocompatibilité). On peut donc comparer ce que dirait un \og single-run \fg{} (juste ce dernier run) avec la classification multi-run.

\begin{center}
\begin{tabular}{@{}lrr@{}}
\toprule
\textbf{Classification multi-run} & \textbf{single=plan} & \textbf{single=non-plan} \\
\midrule
\texttt{always\_planar}      & 362 & 0  \\
\texttt{mostly\_planar}      &   3 & 0  \\
\texttt{unstable}            &  35 & 44 \\
\texttt{mostly\_non\_planar} &   5 & 16 \\
\texttt{always\_non\_planar} &   0 & 54 \\
\texttt{ambiguous}           &   7 & 1  \\
\bottomrule
\end{tabular}
\end{center}

\subsection{Lecture}

\paragraph{Faux positifs single-run.}
Sur les 527 solutions h5, \textbf{40 cas} (\textbf{7,6\%}) seraient classés \og plan \fg{} en single-run alors que la classification multi-run dit \og problématique \fg{} (35 \texttt{unstable}, 5 \texttt{mostly\_non\_planar}). Ces 40 cas sont précisément des bistables où le \emph{dernier} run xTB a eu de la chance et est tombé dans le bassin plan.

\paragraph{Faux négatifs single-run.}
\textbf{0 cas}. Si single-run dit \og non-plan \fg{}, multi-run ne dit jamais \texttt{always\_planar}. \textbf{L'asymétrie est totale.}

\paragraph{Conclusion pratique.}
\begin{itemize}[leftmargin=2em,itemsep=0.2em]
  \item Un single-run \og non-plan \fg{} est \textbf{100\% fiable}~: la molécule \emph{n'est pas} \texttt{always\_planar}.
  \item Un single-run \og plan \fg{} a \textbf{92\% de chances d'être réellement} \texttt{always\_planar} sur h5. Mais il y a 8\% de chances que ce soit en réalité un bistable.
  \item Pour quantifier la fraction de molécules \og toujours planes \fg{} dans un grand jeu de données, le single-run \textbf{surestime de ${\sim}8$ points de pourcentage}.
\end{itemize}

\paragraph{Recommandation.}
Pour les explorations massives (h6, h7), un single-run reste utile pour \emph{exclure} les non-planes (rapidement). Mais pour confirmer qu'une molécule est vraiment \texttt{always\_planar}, il faut le multi-run.

% ======================================================
\section{Effet de la taille $h$}
% ======================================================

\begin{center}
\begin{tabular}{@{}lrrrrr@{}}
\toprule
\textbf{$h$} & \textbf{\# mol.} & \textbf{\# sol.} & \texttt{always\_planar} & \texttt{unstable} & \texttt{always\_non\_planar} \\
\midrule
h3 &  3 &  38 & 100\% &  0\% &  0\% \\
h4 &  6 & 127 &  84\% &  9\% &  5\% \\
h5 & 16 & 527 &  69\% & 15\% & 10\% \\
\bottomrule
\end{tabular}
\end{center}

\subsection{Effet du nombre de défauts}

Quand on regroupe les solutions h5 par couple $(n_5, n_7)$ (nombre de pentagones, nombre d'heptagones)~:

\begin{center}
\begin{tabular}{@{}crrrr@{}}
\toprule
$(n_5, n_7)$ & $n$ & \% AP & \% \texttt{unstable} & \% non-plan \\
\midrule
$(0,0)$ tout-6 & 128 & 100\% &  0\% &  0\% \\
$(1,1)$        & 222 &  60\% & 19\% & 16\% \\
$(2,2)$        & 177 &  56\% & 20\% & 22\% \\
\bottomrule
\end{tabular}
\end{center}

\paragraph{Lecture --- effet de saturation.}
Passer de $0$ défaut à $1$ défaut (= une paire 5-7) fait chuter \% AP de $100\%$ à $60\%$~: chute de \textbf{40 points}. Mais passer de $1$ à $2$ défauts ne fait chuter que de $60\%$ à $56\%$~: \textbf{4 points seulement}. Le coût marginal du second défaut est minime.

C'est une \emph{saturation}~: la première paire 5-7 \og casse \fg{} la situation de planarité parfaite, mais une fois cette cassure faite, ajouter une seconde paire ne dégrade presque pas. Cela suggère que la planarité, une fois compromise localement, ne peut plus être \og pire \fg{} qualitativement~: ce qui change c'est juste la magnitude de l'angle.

\begin{hypothese}
La planarité est un état \og fragile à 1 défaut\fg{}~: c'est l'introduction du \emph{premier} défaut isolé qui décide. Au-delà, ce qui compte est la \emph{position relative} des défauts (sont-ils appariés? alignés sur un axe?), pas leur nombre.
\end{hypothese}

\subsection{Prédictions pour h6+}

Si la tendance se prolonge, h6 devrait avoir~:
\begin{itemize}[leftmargin=2em,itemsep=0.2em]
  \item ${\sim}55\text{--}60\%$ \texttt{always\_planar} global,
  \item mais les molécules à topologie compacte ou linéaire devraient rester \og bonnes \fg{} ($> 85\%$ AP),
  \item et les zigzags ou ramifications devraient s'effondrer ($< 20\%$ AP).
\end{itemize}
C'est testable directement sur les 68 fichiers \texttt{.graph} h6 disponibles dans \texttt{csp\_solver/experiments/plane/benzdb/h6/}.

% ======================================================
\section{Effet des configurations CSP}\label{sec:configs}
% ======================================================

\subsection{Précision vs couverture}

Les 8 configurations CSP combinent 3 flags~: \texttt{--no-freeze}, \texttt{--no-table}, \texttt{--adj-57}. Sur h5~:

\begin{center}
\begin{tabular}{@{}lrrr@{}}
\toprule
\textbf{Config} & \textbf{\# sol.} & \textbf{\# signatures} & \textbf{\% AP} \\
\midrule
\texttt{adj-57}                       &  16 &  1 & \textbf{100\%} \\
\texttt{adj-57\_no-table}             &  45 & 22 & 98\%           \\
\texttt{default}                      &  36 &  9 & 72\%           \\
\texttt{no-freeze}                    &  32 & 11 & 75\%           \\
\texttt{no-table}                     &  90 & 34 & 73\%           \\
\texttt{no-freeze\_no-table}          & 194 & 41 & 51\%           \\
\texttt{adj-57\_no-freeze\_no-table}  &  94 & 31 & 74\%           \\
\bottomrule
\end{tabular}
\end{center}

\paragraph{Lecture en termes de précision/couverture.}
\begin{itemize}[leftmargin=2em,itemsep=0.2em]
  \item \textbf{Précision} = parmi les solutions générées, fraction qui est réellement \texttt{always\_planar}. Reflète à quel point les solutions générées sont \og bonnes \fg{}.
  \item \textbf{Couverture} = nombre de solutions distinctes générées (signatures). Reflète à quel point on explore l'espace.
  \item \texttt{adj-57} a précision parfaite (100\%) mais couverture minimale (1 signature en h5)~: \textbf{on ne génère qu'une seule structure substituée par molécule}.
  \item \texttt{no-freeze\_no-table} a couverture maximale (41 signatures distinctes en h5) mais précision médiocre (51\%)~: une solution sur deux est non-planaire.
\end{itemize}

\paragraph{Stabilité cross-config.}
Une question importante~: si une même molécule + même signature apparaît dans plusieurs configs, donne-t-elle la même classification?

Sur 30 cas mol+sig partagés entre $\geq 2$ configs avec classifications différentes~:
\begin{itemize}[leftmargin=2em,itemsep=0.2em]
  \item Les classes différentes sont \emph{toujours voisines}~: \texttt{unstable} $\leftrightarrow$ \texttt{mostly\_non\_planar}, ou \texttt{mostly\_non\_planar} $\leftrightarrow$ \texttt{always\_non\_planar}.
  \item \textbf{Aucun saut} entre \texttt{always\_planar} et \texttt{always\_non\_planar} dans ces 30 cas.
\end{itemize}

\begin{encadre}{Conclusion structurelle}
La classification est \textbf{largement intrinsèque} au couple (molécule, signature de cycles, topologie réelle). Le parcours du solveur CSP filtre plus ou moins, mais ne crée pas de classifications différentes. Les variations cross-config sont du \emph{bruit aux frontières} (cas borderline qui basculent), pas de vraies inversions.
\end{encadre}

% ======================================================
\section{Cas atypiques candidats à $N=30$}
% ======================================================

\subsection{Distinguer bimodal vs borderline}

Il faut distinguer deux types de \og cas limites \fg{}~:
\begin{itemize}[leftmargin=2em,itemsep=0.3em]
  \item \textbf{Bimodal} (= 91 solutions \texttt{unstable})~: deux groupes nets ($\sim 0\degree$ et $\sim 25\text{--}35\degree$). Ce ne sont \emph{pas} des borderline. $N=30$ ne change pas la classe, juste la probabilité estimée.
  \item \textbf{Borderline} (= ${\sim}21$ solutions identifiées)~: les angles individuels sont \emph{tous} entre $7\degree$ et $13\degree$, autour du seuil de 10\degree. Pas de bimodalité, juste de la dispersion près du seuil. \textbf{C'est ici que $N=30$ aurait du sens}.
\end{itemize}

\subsection{Exemple borderline typique}

\texttt{h4 / 0-6-12-17 [v0=5 v1=7 v2=6 v3=6]}~:
\begin{itemize}[leftmargin=2em,itemsep=0.2em]
  \item Config \texttt{adj-57\_no-freeze\_no-table}~: angles $= [1{,}6;\ 1{,}4;\ 1{,}5;\ 9{,}4;\ 0{,}3;\ 1{,}3;\ 1{,}1;\ 0{,}3;\ 8{,}9;\ 1{,}0]$ $\Rightarrow$ classé \texttt{always\_planar}.
  \item Config \texttt{no-freeze\_no-table}~: angles $= [9{,}5;\ 9{,}9;\ 0{,}2;\ 8{,}8;\ 0{,}9;\ 0{,}9;\ 0{,}5;\ 9{,}2;\ 0{,}7;\ 9{,}6]$ $\Rightarrow$ classé \texttt{ambiguous}.
\end{itemize}
Ici on voit clairement que les angles \og non-plans \fg{} sont $\approx 9\degree$, juste sous le seuil~: pas de second bassin physique, juste une dispersion dont la queue chatouille le seuil. $N=30$ donnerait un classement stable.

\subsection{Liste prioritaire pour re-runs ciblés}

\begin{enumerate}[leftmargin=2em,itemsep=0.3em]
  \item \textbf{21 borderlines} ($\mu \in [7\degree, 13\degree]$ et \emph{non} \texttt{unstable})~: re-run $N=30$. Coût $= 630$ runs. Permet de trancher si vraiment près du seuil. \textbf{Priorité 1}.
  \item \textbf{4 \texttt{always\_non\_planar} avec $\mu < 15\degree$}~: vérifier qu'ils ne sont pas en réalité bistables avec un bassin plan rare. \textbf{Priorité 2}.
  \item \textbf{Multi-runs sur $\sim 20$ molécules h6} ciblées (10 compactes, 10 zigzags)~: tester l'hypothèse topologique sur un $h$ supérieur. \textbf{Priorité 3}.
\end{enumerate}

% ======================================================
\section{Synthèse des hypothèses et incertitudes}
% ======================================================

\subsection{Hypothèses dégagées}

\begin{enumerate}[leftmargin=2em,itemsep=0.3em]
  \item \textbf{Bistabilité géométrique}~: les molécules classées \texttt{unstable} possèdent deux conformères stables (plat et plié) plutôt que d'être \og incertaines numériquement \fg{}.
  \item \textbf{Symétrie de signature}~: une signature palindrome de tailles est associée à un taux d'\texttt{always\_planar} très supérieur (94\% vs 58\% sur h5).
  \item \textbf{Compensation locale 5-7}~: l'appariement local des pentagones et heptagones (chaque 5 a un 7 voisin et vice-versa) augmente fortement la planarité (84\% vs 48\% en h4, 73\% vs 41\% en h5).
  \item \textbf{Topologie parente}~: les molécules à topologie compacte ou linéaire tolèrent les substitutions 5/7, les zigzags et structures branchées les tolèrent mal.
  \item \textbf{Saturation des défauts}~: passer de 0 à 1 paire 5-7 fait chuter la planarité de 40 points; passer de 1 à 2 paires ne fait chuter que de 4 points.
  \item \textbf{Asymétrie single-run}~: un single-run sous-estime systématiquement la non-planarité ($\sim 8\%$ de faux positifs sur h5), jamais l'inverse.
  \item \textbf{Stabilité intrinsèque}~: la classification (mol, sig) est intrinsèque~; les configs CSP filtrent mais ne changent pas la nature géométrique.
\end{enumerate}

\subsection{Incertitudes principales}

\begin{enumerate}[leftmargin=2em,itemsep=0.3em]
  \item \textbf{Statistique faible sur h4} (6 molécules) pour les patterns topologiques.
  \item \textbf{Confusion possible} entre \og adjacence 5-7 \fg{} et \og topologie compacte \fg{}~: les molécules compactes ont plus de chance d'avoir des 5-7 adjacents par construction. Régression jointe nécessaire.
  \item \textbf{Le \og palindrome de signature \fg{}} est un proxy imparfait de la vraie symétrie géométrique (l'ordre des $v_i$ n'est pas canonisé).
  \item \textbf{Le seuil $10\degree$} est cohérent avec la bimodalité globale (creux entre $1\degree$ et $28\degree$), mais arbitraire pour les ${\sim}20$ cas borderline. Un seuil empirique à $5\degree$ serait plus défendable mathématiquement.
  \item \textbf{Toutes les analyses sont sur h3, h4, h5} de BenzAI, donc des benzénoïdes plans simples. Extrapolation à des structures plus exotiques (radicaux, hétéroatomes) non testée.
\end{enumerate}

\subsection{Conclusion}

Le pipeline est cohérent et les patterns sont nets et reproduits sur deux tailles indépendantes ($h=4$ et $h=5$). Les trois mécanismes principaux qui décident de la planarité sont, par ordre d'importance~:
\begin{enumerate}[leftmargin=2em,itemsep=0.2em]
  \item la \textbf{topologie} de la molécule parente (compacte/linéaire vs zigzag/branchée),
  \item l'\textbf{adjacence locale 5-7} (compensation des défauts angulaires en motifs azulène-like),
  \item la \textbf{symétrie} de la distribution des défauts (palindrome).
\end{enumerate}
La contrainte CSP \texttt{--adj-57} encode déjà le mécanisme (2) et garantit 100\% de planarité, au prix d'une couverture très réduite. Pour la suite, le bénéfice marginal le plus important serait de \textbf{coder explicitement la topologie parente} (compacte/linéaire/zigzag) comme prior dans le solveur, et d'introduire une contrainte de symétrie pour augmenter le taux de planarité tout en conservant une couverture intéressante.

\end{document}
